\documentclass[10pt,a4paper]{amsart}
\usepackage{amsmath,amssymb,amsthm,graphicx}
\usepackage{verbatim}
\newenvironment{reference}{\comment}{\endcomment}
%\newenvironment{reference}{}{}
\newenvironment{slogan}{\comment}{\endcomment}
\newenvironment{history}{\comment}{\endcomment}

% For commutative diagrams we use Xy-pic
\usepackage[all]{xy}

% We use 2cell for 2-commutative diagrams.
\xyoption{2cell}
% Theorem environments.
%
\theoremstyle{plain}
\newtheorem{theorem}[subsection]{Theorem}
\newtheorem{proposition}[subsection]{Proposition}
\newtheorem{lemma}[subsection]{Lemma}

\theoremstyle{definition}
\newtheorem{definition}[subsection]{Definition}
\newtheorem{example}[subsection]{Example}
\newtheorem{exercise}[subsection]{Exercise}
\newtheorem{situation}[subsection]{Situation}

\theoremstyle{remark}
\newtheorem{remark}[subsection]{Remark}
\newtheorem{remarks}[subsection]{Remarks}

\numberwithin{equation}{subsection}

% Macros
%
\def\lim{\mathop{\mathrm{lim}}\nolimits}
\def\colim{\mathop{\mathrm{colim}}\nolimits}
\def\Spec{\mathop{\mathrm{Spec}}}
\def\Hom{\mathop{\mathrm{Hom}}\nolimits}
\def\Ext{\mathop{\mathrm{Ext}}\nolimits}
\def\SheafHom{\mathop{\mathcal{H}\!\mathit{om}}\nolimits}
\def\SheafExt{\mathop{\mathcal{E}\!\mathit{xt}}\nolimits}
\def\Sch{\mathit{Sch}}
\def\Mor{\mathop{\mathrm{Mor}}\nolimits}
\def\Ob{\mathop{\mathrm{Ob}}\nolimits}
\def\Sh{\mathop{\mathit{Sh}}\nolimits}
\def\NL{\mathop{N\!L}\nolimits}
\def\CH{\mathop{\mathrm{CH}}\nolimits}
\def\proetale{{pro\text{-}\acute{e}tale}}
\def\etale{{\acute{e}tale}}
\def\QCoh{\mathit{QCoh}}
\def\Ker{\mathop{\mathrm{Ker}}}
\def\Im{\mathop{\mathrm{Im}}}
\def\Coker{\mathop{\mathrm{Coker}}}
\def\Coim{\mathop{\mathrm{Coim}}}

% Boxtimes
%
\DeclareMathSymbol{\boxtimes}{\mathbin}{AMSa}{"02}

%
% Macros for moduli stacks/spaces
%
\def\QCohstack{\mathcal{QC}\!\mathit{oh}}
\def\Cohstack{\mathcal{C}\!\mathit{oh}}
\def\Spacesstack{\mathcal{S}\!\mathit{paces}}
\def\Quotfunctor{\mathrm{Quot}}
\def\Hilbfunctor{\mathrm{Hilb}}
\def\Curvesstack{\mathcal{C}\!\mathit{urves}}
\def\Polarizedstack{\mathcal{P}\!\mathit{olarized}}
\def\Complexesstack{\mathcal{C}\!\mathit{omplexes}}
% \Pic is the operator that assigns to X its picard group, usage \Pic(X)
% \Picardstack_{X/B} denotes the Picard stack of X over B
% \Picardfunctor_{X/B} denotes the Picard functor of X over B
\def\Pic{\mathop{\mathrm{Pic}}\nolimits}
\def\Picardstack{\mathcal{P}\!\mathit{ic}}
\def\Picardfunctor{\mathrm{Pic}}
\def\Deformationcategory{\mathcal{D}\!\mathit{ef}}
\usepackage{cleveref}
\crefname{originallemma}{original lemma}{original lemmas}
\Crefname{originallemma}{Original lemma}{Original lemmas}
\makeatletter
\newlabel{flat-lemma-finite-type-flat-along-fibre-free-variant}{{05L1}{source}{}{sourcepointer-1}{}}
\newlabel{flat-lemma-finite-type-flat-along-fibre-free-variant@cref}{{[originallemma][0][]05L1}{[1][0][]source}{}{}{}}
\makeatother
\crefname{originallemma}{original lemma}{original lemmas}
\Crefname{originallemma}{Original lemma}{Original lemmas}
\makeatletter
\newlabel{spaces-morphisms-lemma-base-change-module-flat}{{05VW}{source}{}{sourcepointer-2}{}}
\newlabel{spaces-morphisms-lemma-base-change-module-flat@cref}{{[originallemma][0][]05VW}{[1][0][]source}{}{}{}}
\makeatother
\title{Etale local freeness along the fibre for a finite type sheaf flat over the base}
\author{The Stacks Project Authors}
\date{stacks-project:spaces-flat.tex:local-snapshot-20261012; GNU FDL 1.2 or later}
\begin{document}
\maketitle
\noindent\textit{Editorial note.} All delivered mathematical text is the author's own, copied from the pinned Stacks Project chapter file spaces-flat.tex; no statement or proof step was rewritten, reordered or removed.\par\medskip
\noindent\textit{Editorial note.} The counted proof is the author's complete argument: a two step etale descent construction: affine etale localization of the base followed by an affine surjective etale cover, flatness transferred by base change, the scheme case lemma, and composition of the steps\par\medskip
\noindent\textit{Editorial note.} The label lemma-finite-type-flat-along-fibre-free-variant names the counted statement, and the proof environment immediately following it in the source is delivered verbatim as its proof.\par\medskip
\noindent\textit{Editorial note.} The counted statement is carried by the author's own lemma environment in the chapter source.\par\medskip
\noindent\textit{Allowed standard prerequisites (editorial).}\begin{itemize}
\item Base change of flatness of a module along morphisms of algebraic spaces
\item Existence of affine etale neighbourhoods of a point and of affine schemes surjective etale over a quasi-compact algebraic space
\item The scheme case: the freeness variant of the finite type flat along the fibre lemma for schemes
\end{itemize}\medskip
\ifdefined\hypertarget\hypertarget{sourcepointer-1}{}\fi
\noindent\textit{Source reference.} Original lemma 05L1: sources/stacks-chow/tags.txt, source lines 7220--7220. Prerequisite pointer invoked at proof line(s) 819; it names a standard fact of a sibling chapter that the counted argument uses, and every step of the counted proof built on it is retained verbatim in the delivered body.\par\medskip
\ifdefined\hypertarget\hypertarget{sourcepointer-2}{}\fi
\noindent\textit{Source reference.} Original lemma 05VW: sources/stacks-chow/tags.txt, source lines 7565--7565. Prerequisite pointer invoked at proof line(s) 816; it names a standard fact of a sibling chapter that the counted argument uses, and every step of the counted proof built on it is retained verbatim in the delivered body.\par\medskip
\begin{lemma}
\label{lemma-finite-type-flat-along-fibre-free-variant}
Let $S$ be a scheme. Let $f : X \to Y$ be a morphism of algebraic spaces
over $S$. Let $\mathcal{F}$ be a quasi-coherent sheaf on $X$.
Let $y \in |Y|$. Set $F = f^{-1}(\{y\}) \subset |X|$. Assume that
\begin{enumerate}
\item $f$ is of finite type,
\item $\mathcal{F}$ is of finite type, and
\item $\mathcal{F}$ is flat over $Y$ at all $x \in F$.
\end{enumerate}
Then there exists an \'etale morphism $(Y', y') \to (Y, y)$
where $Y'$ is a scheme and a commutative diagram of algebraic spaces
$$
\xymatrix{
X \ar[d] & X' \ar[l]^g \ar[d] \\
Y & \Spec(\mathcal{O}_{Y', y'}) \ar[l]
}
$$
such that $X' \to X \times_Y \Spec(\mathcal{O}_{Y', y'})$
is \'etale, $|X'_{y'}| \to F$ is surjective, $X'$ is affine,
and $\Gamma(X', g^*\mathcal{F})$ is a free $\mathcal{O}_{Y', y'}$-module.
\end{lemma}

\begin{proof}
Choose an \'etale morphism $(Y', y') \to (Y, y)$ where $Y'$ is an
affine scheme. Then $X \times_Y Y'$ is quasi-compact.
Choose an affine scheme $X'$ and a surjective \'etale morphism
$X' \to X \times_Y Y'$. Picture
$$
\xymatrix{
X \ar[d] & X' \ar[l]^g \ar[d] \\
Y & Y' \ar[l]
}
$$
Then $\mathcal{F}' = g^*\mathcal{F}$ is flat over $Y'$ at all
points of $X'_{y'}$, see Morphisms of Spaces, Lemma
\ref{spaces-morphisms-lemma-base-change-module-flat}.
Hence we can apply the lemma in the case of schemes
(More on Flatness, Lemma
\ref{flat-lemma-finite-type-flat-along-fibre-free-variant})
to the morphism
$X' \to Y'$, the quasi-coherent sheaf $g^*\mathcal{F}$, and the point $y'$.
This gives an \'etale morphism $(Y'', y'') \to (Y', y')$ and a commutative
diagram
$$
\xymatrix{
X \ar[d] & X' \ar[l]^g \ar[d] & X'' \ar[l]^{g'} \ar[d] \\
Y & Y' \ar[l] & \Spec(\mathcal{O}_{Y'', y''}) \ar[l]
}
$$
To get what we want we take $(Y'', y'') \to (Y, y)$
and $g \circ g' : X'' \to X$.
\end{proof}

\end{document}
